\section{The zero divisors}\label{sec:factors}

\begin{proof}[Proof of \Cref{thm:main}]
Use the graphs of \Cref{prop:no-arrangements} and the complex \(X\) constructed in \Cref{sec:topology}. It is finite and two-dimensional. By \Cref{cor:torsion-free}, it is a classifying space for the torsion-free group \(G=\pi_1(X)\), which is finitely presented.

Let \(A'\) be the full vertex set of the component of \(x_A\) and \(B'\) that of the component of \(x_B\). For \(x\in A'\), let \(g_x\in G\) be the label of a path from \(x_A\) to \(x\); for \(y\in B'\), define \(h_y\) using paths from \(x_B\). These labels do not depend on the path, because labels of all closed walks in each component are killed by its cone. We multiply in path order.

Set
\[
\alpha=\sum_{x\in A'}g_x,\qquad
\beta=\sum_{y\in B'}h_y^{-1}
\quad\text{in }\mathbb F_2[G].
\]
Both sums are finite. The root contributes the identity to each. By \Cref{prop:topology}, no other vertex contributes identity, so its coefficient is one in each factor. Thus \(\alpha\ne0\) and \(\beta\ne0\). No assertion about equality between two non-root labels is needed.

On \(A'\times B'\), form the graph of simultaneous steps of a common signed label. A step
\[
x\xrightarrow{t}x',\qquad y\xrightarrow{t}y'
\]
sends \(g_x\) to \(g_xt\) and \(h_y\) to \(h_yt\), and therefore preserves
\[
g_{x'}h_{y'}^{-1}
 =(g_xt)(h_yt)^{-1}
 =g_xh_y^{-1}.
\]
The reverse step has label \(\bar t\). Since full components are used, the degree at \((x,y)\) is the entire number \(|S_x\cap S_y|\), which is odd by \eqref{eq:odd-intersections}. The girth condition excludes loops for the chosen large \(n\); parallel edges would not affect the degree-sum argument.

Every finite connected component of the simultaneous-step graph has an even number of vertices, since the sum of its odd degrees is twice its number of edges. The value \(g_xh_y^{-1}\) is constant on that component, so the component contributes zero to the expansion of \(\alpha\beta\) over \(\mathbb F_2\). Summing over all components gives \(\alpha\beta=0\). This proves \Cref{thm:main} and disproves the universal zero-divisor assertion.
\end{proof}

In particular, \(G\) does not have the unique-product property. The construction
is probabilistic: it proves that suitable finite matchings exist, without
specifying matchings from which a concrete presentation and zero-divisor
factors can be read off.
